\documentclass[10pt,aspectratio=169]{beamer}

% --- Minimal theme ---------------------------------------------------------
\usetheme{default}
\usecolortheme{dove}
\usefonttheme{professionalfonts}
\setbeamertemplate{navigation symbols}{}
\setbeamertemplate{footline}[frame number]
\setbeamertemplate{itemize items}[circle]
\setbeamertemplate{section in toc}[sections numbered]
\setbeamerfont{frametitle}{size=\large,series=\bfseries}
\setbeamercolor{frametitle}{fg=black}
\setbeamercolor{title}{fg=black}
\setbeamercolor{structure}{fg=black}

\usepackage{booktabs}
\usepackage{array}
\usepackage{graphicx}
\graphicspath{{../figs/}}

\title{Where do the temperature datasets disagree?}
\subtitle{Country- and time-level diagnostics for five reanalyses}
\author{}
\date{\today}

\begin{document}

\begin{frame}\titlepage\end{frame}

\begin{frame}{Outline}\tableofcontents\end{frame}

% ===========================================================================
\section{Overall agreement}

\begin{frame}{Setup}
  \begin{itemize}
    \item Balanced sample: 4{,}592 (iso, year) pairs observed in all five
          products. 158 countries $\times$ 31 years (1980--2013).
    \item Variable: pop-weighted annual mean temperature (°C).
  \end{itemize}
\end{frame}

\begin{frame}{Raw pairwise disagreement is small --- but misleading}
  \centering
  \includegraphics[height=0.6\textheight]{tc_pairwise_mad.pdf}

  \vspace{0.2em}
  \footnotesize
  Raw mean |diff| is only 0.5--0.7 °C across all 10 pairs, and raw
  correlations are all $\geq 0.992$. That number hides everything: it is
  dominated by the cross-country climate signal, which the datasets agree on
  by construction.
\end{frame}

\begin{frame}{What matters is the \emph{residualised} correlation}
  After partialling out the full FE structure used in the growth regression
  --- country FE, year FE, country-specific linear and quadratic time trends
  --- only the weather variation that identifies the panel coefficient is
  left.
  \vspace{0.4em}

  \centering
  \includegraphics[height=0.58\textheight]{tc_cor_heatmap_resid.pdf}

  \vspace{0.2em}
  \footnotesize
  The range collapses from $[0.992, 0.996]$ raw to $[0.74, 0.96]$ residualised.
  Modern reanalyses (ERA5, JRA-3Q) correlate 0.96 on weather signal.
  GMFD vs.\ MERRA2 drops to 0.74.
\end{frame}

\begin{frame}{Bias depends on temperature level}
  \centering
  \includegraphics[width=0.98\linewidth]{tc_bias_by_level.pdf}

  \vspace{0.2em}
  \footnotesize
  Loess of $A-B$ against $B$. Non-flat curves are the rule, not the exception.
  The cold tail and $\sim\!15$°C band are where most pairs diverge.
\end{frame}

% ===========================================================================
\section{Country-level disagreement}

\begin{frame}{ERA5-025 vs MERRA2: top 20 countries}
  \centering
  \includegraphics[width=0.7\linewidth]{tc_country_bars.pdf}

  \vspace{0.2em}
  \footnotesize
  Costa Rica tops the list at nearly 2\,°C mean |diff|.
  Central American and Andean LAC countries dominate, along with mountainous
  EECA countries (KGZ, TJK).
\end{frame}

\begin{frame}{Most of the country disagreement is a \emph{static bias}}
  \small
  Top 5 ERA5 vs MERRA2 disagreeing countries, decomposed:

  \vspace{0.4em}
  \centering
  \begin{tabular}{l r r r}
    \toprule
    iso & mean bias (°C) & std.\ dev.\ of diff & interpretation \\
    \midrule
    CRI & $-1.92$ & 0.26 & \emph{constant} offset \\
    HTI & $-1.53$ & 0.24 & constant offset \\
    KGZ & $-1.49$ & 0.22 & constant offset \\
    GTM & $-1.45$ & 0.64 & offset + year-to-year variation \\
    KOR & $+1.40$ & 0.15 & constant offset \\
    \bottomrule
  \end{tabular}

  \vspace{0.6em}
  \footnotesize
  When the SD is tiny compared to the bias, the two products disagree about
  the country's \emph{climate}, not about year-to-year \emph{weather}. This
  is consistent with differences in grid resolution, topographic handling,
  and population weighting --- \emph{not} from the climate signal itself.
\end{frame}

% ===========================================================================
\section{Residualised: where do they disagree on weather?}

\begin{frame}{Partial out the full regression FE structure}
  \[
    \texttt{temp}_{ct} = \alpha_c + \tau_t + \delta_c\,t + \phi_c\,t^2 + \tilde T_{ct}
  \]
  Country FE, year FE, and country-specific linear and quadratic time trends
  --- exactly the FE structure of the main growth regression. What remains,
  $\tilde T_{ct}$, is the year-to-year weather variation that identifies the
  panel coefficient. If the datasets agree on this variation,
  $\tilde T^{A}_{ct} \approx \tilde T^{B}_{ct}$.
\end{frame}

\begin{frame}{Residual ERA5 vs MERRA2}
  \centering
  \includegraphics[height=0.72\textheight]{tc_resid_scatter.pdf}

  \vspace{0.2em}
  \footnotesize
  Much tighter than the raw scatter --- weather signals agree better than
  climate levels --- but noticeable spread remains in the tails.
\end{frame}

\begin{frame}{Top countries by \emph{residual} disagreement}
  \centering
  \includegraphics[width=0.7\linewidth]{tc_resid_country_bars.pdf}

  \vspace{0.2em}
  \footnotesize
  Removing climate reorders the list: tropical and semi-arid countries
  (CAF, RWA, BDI, HND, AGO, GTM) now dominate. These are where datasets
  disagree about year-to-year \emph{weather}, not just about the climate
  mean --- exactly the variation our panel regression uses.
\end{frame}

\begin{frame}{Residual temperature over time: top 8 vs bottom 8}
  \centering
  \includegraphics[height=0.82\textheight]{tc_resid_ts_top.pdf}\hfill
  \includegraphics[height=0.82\textheight]{tc_resid_ts_bottom.pdf}
\end{frame}

\begin{frame}{Residual time series: what to notice}
  \begin{itemize}
    \item Left: the 8 countries where the five datasets disagree most on
          residual weather. Right: the 8 where they agree most.
    \item In the \textbf{top-8 (disagreement) panel} the five lines fan out
          by 0.5--1.0\,°C in individual years, and the \emph{sign} of the
          anomaly can flip across datasets. UDel (station-interpolation,
          not a reanalysis) is often the outlier.
    \item In the \textbf{bottom-8 (agreement) panel} the five lines nearly
          coincide year-by-year --- the reanalyses agree on weather.
    \item The two sets are geographically distinct: disagreement is
          concentrated in tropical / mountainous countries; agreement in
          large, well-observed temperate countries.
  \end{itemize}
\end{frame}

% ===========================================================================
\section{Disagreement, income, and climate}

\begin{frame}{Where in the world is disagreement concentrated?}
  For each country, average residualised MAD across all 10 pairs
  $\Rightarrow$ a single scalar ``how much do the five datasets disagree
  about this country's weather'' per iso.

  \vspace{0.3em}

  Then cross with country-mean log GDP per capita (PPP) and country-mean
  temperature.
  \vspace{0.6em}

  \centering
  \begin{tabular}{l r r}
    \toprule
    & \textbf{Residualised MAD} & Raw MAD \\
    \midrule
    Corr.\ with log GDP per cap & $\mathbf{-0.45}$ & $-0.01$ \\
    Corr.\ with mean temperature & $\mathbf{+0.33}$ & $+0.08$ \\
    \bottomrule
  \end{tabular}

  \vspace{0.6em}
  \footnotesize
  Raw disagreement is uncorrelated with development or climate --- the
  static biases are scattered without a clear pattern. But the \emph{weather
  signal} disagreement is strongly concentrated in poor, hot countries.
\end{frame}

\begin{frame}{Disagreement vs income}
  \centering
  \includegraphics[width=0.85\linewidth]{tc_disagr_vs_income.pdf}

  \vspace{0.2em}
  \footnotesize
  Each point is a country; $y$ is the average residualised pairwise MAD
  across 10 pairs. Dashed = OLS. Richer countries are measured consistently
  across reanalyses; in the poorest countries the five products disagree
  by up to 0.4\,°C per year \emph{after} removing climate and trend.
\end{frame}

\begin{frame}{Disagreement vs climate}
  \centering
  \includegraphics[width=0.85\linewidth]{tc_disagr_vs_temp.pdf}

  \vspace{0.2em}
  \footnotesize
  OLS fit (dashed). Hot countries show more weather-signal disagreement.
  Combined with the income plot, the worst-measured countries are hot and
  poor --- the tropical low-income band.
\end{frame}

\begin{frame}{For contrast: raw MAD is flat in both dimensions}
  \centering
  \includegraphics[width=0.48\linewidth]{tc_disagr_raw_vs_income.pdf}
  \includegraphics[width=0.48\linewidth]{tc_disagr_raw_vs_temp.pdf}

  \vspace{0.2em}
  \footnotesize
  Static bias (raw |diff|) bears no systematic relationship to income or
  climate --- it's driven by idiosyncratic grid/topography effects. It's
  only the \emph{residualised} disagreement that sorts along development
  and climate.
\end{frame}

% ===========================================================================
\section{Trends and outliers}

\begin{frame}{Per-country warming trends: ERA5 vs MERRA2}
  \centering
  \includegraphics[height=0.7\textheight]{tc_trend_scatter.pdf}

  \vspace{0.2em}
  \footnotesize
  Scatter of country-level slope coefficients from OLS of temp on year.
  Trends line up tightly on the 45° line --- the datasets agree on
  \emph{how fast} each country is warming, even where they disagree on the
  level.
\end{frame}

\begin{frame}{Temporal evolution of disagreement}
  \centering
  \includegraphics[width=\linewidth]{tc_year_series.pdf}

  \vspace{0.2em}
  \footnotesize
  Mostly flat --- \emph{except} for a visible bump in 2009--2010 for
  several pairs. Worth drilling in.
\end{frame}

\begin{frame}{Zoom 2005--2013: the spike is a GMFD event}
  \centering
  \includegraphics[width=0.95\linewidth]{tc_year_zoom.pdf}

  \vspace{0.2em}
  \footnotesize
  GMFD pairs (solid lines) all jump in 2010; pairs not involving GMFD
  (dotted) are flat. Only four datasets --- not all five --- have a 2010
  anomaly.
\end{frame}

\begin{frame}{GMFD vs the ensemble of the other four}
  \centering
  \includegraphics[width=0.9\linewidth]{tc_gmfd_year.pdf}

  \vspace{0.2em}
  \small
  Mean |GMFD $-$ (ERA5 + JRA-3Q + MERRA2 + UDel)/4| over countries, by year.
  2010 is the largest year in the sample (0.70\,°C, $+35\%$ above the
  long-run mean); 2009 is the second-largest. No comparable spike occurs
  elsewhere.
\end{frame}

\begin{frame}{Where does the 2010 deviation land?}
  \centering
  \includegraphics[height=0.75\textheight]{tc_2010_countries.pdf}

  \vspace{0.1em}
  \footnotesize
  Signed GMFD $-$ mean-of-others for 2010. Spatially \textbf{coherent}:
  GMFD is $+1.6$ to $+2.5$\,°C warmer in Northern Europe (NOR, SWE, IRL,
  DNK, BEL, FRA); $-1.6$ to $-2.3$\,°C cooler across the Middle East
  (IRQ, YEM, ARM, GEO, OMN, SAU, JOR, KWT).
\end{frame}

\begin{frame}{2009 shows the same signature, smaller amplitude}
  \centering
  \includegraphics[height=0.75\textheight]{tc_2009_countries.pdf}

  \vspace{0.1em}
  \footnotesize
  2009 pre-echoes 2010: same Northern-Europe-warm / Middle-East-cool
  pattern, fewer countries and smaller magnitude. Suggests a progressive
  input change rather than a one-year mistake.
\end{frame}

\begin{frame}{What's going on? Probably a GMFD input transition}
  \begin{itemize}
    \item GMFD (Princeton Global Meteorological Forcing Dataset, v3) blends
          multiple underlying products; its newer years rely on different
          observational inputs than its historical record (CRU TS is the
          classic backbone and has its own post-2010 quality issues).
    \item A spatially coherent bias of $\pm 2$\,°C at continent scale
          is hard to explain as weather: the 2010 Russian heat wave and
          Pakistan floods don't project onto ``Northern Europe warm,
          Middle East cool''.
    \item Most likely \textbf{an input/calibration change in GMFD around
          2009--2010} that the reanalyses do not share.
    \item \textbf{Implication}: if you use GMFD for panel damage work, the
          2010 cell contains a non-climate shock. Drop or flag these years,
          or use one of the reanalyses (ERA5 / JRA-3Q / MERRA2).
  \end{itemize}
\end{frame}

\begin{frame}{Top 20 outlier (iso, year) pairs}
  \centering
  \scriptsize
  \begin{tabular}{l r r r r}
    \toprule
    iso & year & ERA5-025 & MERRA2 & diff \\
    \midrule
    GTM & 1986 & 19.5 & 22.2 & $-2.69$ \\
    GTM & 1985 & 19.5 & 22.2 & $-2.66$ \\
    RWA & 1999 & 17.4 & 20.0 & $-2.62$ \\
    GTM & 1987 & 20.1 & 22.7 & $-2.58$ \\
    BDI & 1999 & 18.8 & 21.4 & $-2.55$ \\
    HND & 1994 & 22.6 & 25.0 & $-2.43$ \\
    CRI & 1996 & 21.6 & 24.0 & $-2.38$ \\
    HND & 1987 & 22.5 & 24.8 & $-2.38$ \\
    RWA & 1992 & 17.6 & 20.0 & $-2.36$ \\
    CRI & 1993 & 21.9 & 24.2 & $-2.28$ \\
    \bottomrule
  \end{tabular}

  \vspace{0.4em}
  \footnotesize
  Top extremes are concentrated in a handful of Central American and East
  African countries. MERRA2 is consistently warmer for the same country-year.
\end{frame}

% ===========================================================================
\section{Takeaways}

\begin{frame}{What we learned}
  \begin{itemize}
    \item Raw agreement is high ($r \geq 0.992$), but this is dominated by
          the cross-country climate signal. On \textbf{residualised weather
          variation} pairwise correlations range from 0.74 (GMFD vs MERRA2)
          to 0.96 (ERA5 vs JRA-3Q) --- a much more honest picture of where
          the panel ID disagrees.
    \item \textbf{Country-level disagreement is largely a static bias} in
          ~30 countries --- mostly mountainous, coastal, or small:
          Costa Rica, Haiti, Kyrgyzstan, Guatemala, Korea. Standard deviations
          of the bias are tiny compared to its level.
    \item \textbf{Weather-signal disagreement} (residualised) is
          concentrated in the tropics: CAF, RWA, BDI, HND, AGO.
          This is the disagreement that \emph{enters our panel regression} and
          plausibly drives the ME gap from the previous deck.
    \item Trends agree well --- all datasets see roughly the same warming
          per country.
    \item For panel damage work, the \emph{residualised} disagreement is what
          matters, and it is spatially non-random.
    \item \textbf{It sorts along development and climate:} correlation with
          log GDP pc $= -0.45$; correlation with mean temperature $= +0.33$.
          The datasets are most ambiguous precisely for the countries most
          exposed to climate damage.
  \end{itemize}
\end{frame}

\end{document}
